%! Logarithmic Function Context and Data Modeling — Unit 2, Lesson 2.14 — Homework (Answer Key)
\documentclass[10pt]{article}
\usepackage{precalculus-article}
\usepackage{precalculus-key}   % pulls in -boxes; do NOT also load -boxes
\newcommand{\TallMath}[1]{$\displaystyle #1\rule[-1.4em]{0pt}{3.2em}$}

\begin{document}

\pageheader{Unit 2, Lesson 2.14}{Homework --- Answer Key}
\namedateperiod

\vspace{0.1in}

\begin{enumerate}[leftmargin=1.8em, itemsep=14pt]

\item Classify each table.

\vspace{6pt}
\begin{center}
\renewcommand{\arraystretch}{1.8}
\begin{tabular}{cc}
  \begin{tabular}{|c|c|}
  \hline
  \rowcolor{sky}
  $x$ & $y$ \\
  \hline
  $1$ & $0$ \\
  \hline
  $3$ & $1$ \\
  \hline
  $9$ & $2$ \\
  \hline
  $27$ & $3$ \\
  \hline
  \end{tabular}
  \qquad\qquad &
  \begin{tabular}{|c|c|}
  \hline
  \rowcolor{sky}
  $x$ & $y$ \\
  \hline
  $0$ & $2$ \\
  \hline
  $1$ & $6$ \\
  \hline
  $2$ & $18$ \\
  \hline
  $3$ & $54$ \\
  \hline
  \end{tabular}
\end{tabular}
\end{center}

\vspace{4pt}
\textbf{Table A:} \quad \ans{Logarithmic} \quad because
\ansline{as $y$ increases by 1, $x$ multiplies by 3 --- proportional input change for
equal output increments is the logarithmic pattern.}

\vspace{4pt}
\textbf{Table B:} \quad \ans{Exponential} \quad because
\ansline{over equal $x$-intervals (add 1), the output multiplies by 3 --- proportional
output change for equal input increments is the exponential pattern.}

\item Model from $f(1)=0$ and $f(e^3)=6$, $f(x)=a\ln(x)$.

\begin{enumerate}[leftmargin=1.5em, itemsep=8pt, label=\alph*.]
  \item $a\ln(e^3) = 6 \;\Rightarrow\; a \cdot $ \ans{3} $= 6 \;\Rightarrow\; a = $ \ans{2}

  \item $f(x) = $ \ans{$2\ln(x)$}

  \item $f(e^5) = $ \ans{$2\ln(e^5) = 2 \cdot 5 = 10$}
\end{enumerate}

\item Species-count model $S(A) = a\log_{10}(A)+k$.

\begin{enumerate}[leftmargin=1.5em, itemsep=8pt, label=\alph*.]
  \item Equation 1: \ans{$3 = a\log_{10}(1)+k = k$} \quad Equation 2: \ans{$6 = a\log_{10}(10)+k = a+k$}

  \item $k = $ \ans{3} \qquad $a = $ \ans{3}

  \item $S(A) = $ \ans{$3\log_{10}(A)+3$} \qquad $S(10000) = $ \ans{$3(4)+3 = 15$}
\end{enumerate}

\item $f(x) = 2\log_2(x-1)+3$.

\begin{enumerate}[leftmargin=1.5em, itemsep=8pt, label=\alph*.]
  \item Domain: \ans{$x > 1$, i.e.\ $(1,\,\infty)$}

  \item $f(5) = 2\log_2(5-1)+3 = 2\log_2(\ans{4})+3 = 2\cdot\ans{2}+3 = $ \ans{7}

  \item Solve $f(x)=7$:
  \ansline{$2\log_2(x-1)+3 = 7 \;\Rightarrow\; 2\log_2(x-1) = 4 \;\Rightarrow\;
  \log_2(x-1) = 2 \;\Rightarrow\; x-1 = 4 \;\Rightarrow\; x = 5$}
  $x = $ \ans{5}
\end{enumerate}

\item AP-style MC: \ans{(B)} --- A child's height grows rapidly then slows, which is
      best modeled by a logarithmic function (increasing but concave down, decelerating
      growth). Choice (A) and (D) are exponential; (C) is linear.

\item Extension:
\ansline{By change-of-base, $\log_{10}(x) = \dfrac{\ln x}{\ln 10}$, so
$a\log_{10}(x) = \dfrac{a}{\ln 10}\cdot\ln x$. This means $g(x) = a\log_{10}(x)$
is a vertical dilation of $f(x) = a\ln(x)$ by the factor $\dfrac{1}{\ln 10} \approx 0.4343$.
Equivalently, $f(x) = \ln(10)\cdot g(x)$, so $f$ is a vertical stretch of $g$ by
$\ln 10 \approx 2.303$.}

\end{enumerate}

\vspace{0.12in}

\begin{reflectionbox}
What is one thing from today's lesson that you feel confident about, and one question you
still have about constructing or using logarithmic models?

\writelines{2}
\end{reflectionbox}

\end{document}
